% ===========================================================================
% Slides preamble (beamer target) — loaded by drivers/slides.tex.
% Extracted from the donor course (PHYS 567 shared/preamble.tex).
% Builds run from the REPO ROOT: all paths are root-relative.
% ===========================================================================
\usepackage[T1]{fontenc}
\usepackage{lmodern}
\usepackage{amsmath,amssymb,mathtools}
\usepackage{array}
\usepackage{booktabs}
\usepackage{adjustbox}
\usepackage{tikz}
\usetikzlibrary{arrows.meta,calc,positioning,shapes.geometric}
\usetikzlibrary{overlay-beamer-styles} % visible on=<...> for staged figure reveals

% --- Kit palette (keep in sync with preamble-web.tex and style.css) ----------
\definecolor{physblue}{HTML}{13294B}
\definecolor{physorange}{HTML}{C84113}
\definecolor{physgray}{HTML}{F3F4F6}
\definecolor{physink}{HTML}{111827}
\definecolor{physgreen}{HTML}{1B7A3D}

\setbeamertemplate{navigation symbols}{}
\setbeamertemplate{itemize item}{\raise0.2ex\hbox{\small$\blacktriangleright$}}
\setbeamertemplate{itemize subitem}{\raise0.2ex\hbox{\scriptsize$\bullet$}}
\setbeamertemplate{enumerate item}{\insertenumlabel.}

\setbeamercolor{normal text}{fg=physink,bg=white}
\setbeamercolor{frametitle}{fg=white,bg=physblue}
\setbeamercolor{title}{fg=white,bg=physblue}
\setbeamercolor{structure}{fg=physblue}
\setbeamercolor{block title}{fg=white,bg=physblue}
\setbeamercolor{block body}{fg=physink,bg=physgray}
\setbeamercolor{alerted text}{fg=physorange}

\setbeamerfont{title}{size=\Large,series=\bfseries}
\setbeamerfont{frametitle}{size=\large,series=\bfseries}
\setbeamerfont{block title}{series=\bfseries}

\hypersetup{colorlinks=true, urlcolor=physblue, linkcolor=physblue}

% --- Lecture title metadata --------------------------------------------------
% Content files open with \kitlecturetitle{Topic Title}; the driver has already
% defined \kitlecture (the two-digit lecture number).
\newcommand{\kitlecturetitle}[1]{%
  \title{#1}\subtitle{Lecture \kitlecture}\author{\coursecode}\date{\courseterm}}

% --- \fig{path}{alt-key}: figure with alt text born alongside it -------------
% #1 = root-relative TikZ source path (no .tex), #2 = alt key "lectureNN/name".
% The slides target renders the figure; the alt key is consumed by the web
% target (alt/lectureNN.tex sidecars). Same signature in both targets.
\newcommand{\fig}[2]{\begin{center}\input{#1}\end{center}}

% --- shared-figure alt borrowing (web-semantic; no-op on slides) -------------
% A lecture reusing another lecture's figure declares \usealtfrom{lectureMM}
% near the top so the web target can resolve the owner's alt keys.
\newcommand{\usealtfrom}[1]{}

% --- \case{N}: circled case/piece label --------------------------------------
% Slides render a true circled glyph; the web target renders "(N)" (its
% converter mangles decorated \textcircled arguments — donor lesson 2026-08-02).
\newcommand{\case}[1]{\textcircled{\raisebox{-0.5pt}{\scriptsize #1}}}

% --- Announcements sidecars --------------------------------------------------
% Semester-ephemeral announcements live in announcements/lectureNN.tex, NOT in
% lecture bodies: edit the sidecar the day of lecture, rebuild, done.  An
% absent sidecar renders nothing.  The web target no-ops both macros.
\newcommand{\kitannouncementspath}{announcements/lecture\kitlecture.tex}
\newcommand{\announcementsframe}{%
  \IfFileExists{\kitannouncementspath}%
    {\begin{frame}{Announcements}\input{\kitannouncementspath}\end{frame}}{}}
\newcommand{\announcementsblock}{%
  \IfFileExists{\kitannouncementspath}%
    {\begin{block}{Announcements}\input{\kitannouncementspath}\end{block}}{}}
